尽管本文结果已经清晰证明了求解器层面的改进效果，仍有若干局限性需要如实说明。首先，本文所有长航时 endurance 评测均在 Isaac Sim 的平坦地形上完成。这样的设置非常适合隔离 branch-consistency 失效模式，并验证基于 IPM 的修复在受控环境中的真实效果，但它并未覆盖野外机器人将面临的全部接触现象。特别是，不规则地形与柔性地形会引入更大的接触不确定性、持续变化的支撑几何以及更强的模型失配，这些因素都可能与优化器发生比平地实验更复杂的耦合。因此，一个自然且必要的下一步，是将本文的 IPM + branch-constraint 框架扩展到感知驱动的非平地与柔顺地形运动控制中，以检验这种可行域塑形策略在复杂接触反馈下是否仍然稳健。

其次，尽管本文方法在长航时评测中保持了稳定前行，60~s endurance 下的前向位移并未与时间呈现完全线性增长。0.638~m 的最终位移已经证明系统能够持续前进并维持实时稳定，但它也说明：在极长时域内，控制器会逐渐进入一种更保守的运行状态。我们认为，这一现象更可能来自启发式落足点规划、moving-goal reference 生成以及微小状态观测误差的长期累积耦合，而不是 branch-consistency 约束本身的再次失效。换言之，当“错误分支”这一灾难性病灶被去除后，更细微的上层运动生成问题开始成为限制长期推进效率的主导因素。解决这一非线性位移缩放问题，需要进一步优化轨迹生成器、落足策略和航向调节逻辑，而非继续修改分支一致性约束本身。

这些局限性也反过来说明了本文贡献的边界与价值所在。本文最强的贡献位于“求解器正确性”和“机构可行性建模”层面：我们证明了，物理分支有效性必须被视为优化域的结构属性，而不能留给 runtime 事后修补。未来工作应建立在这个已经修正的求解器基础之上，沿两个方向继续推进：一类是环境复杂度的提升，包括不规则地形、柔性地形和部分可观测地形；另一类是长时域性能塑形，包括更优的落足启发式、更强的 reference 生成机制以及感知、运动意图与全身动力学之间更紧密的耦合。只有同时推动这两个方向，我们才能进一步验证：本文提出的 branch-safe IPM 框架，能否从平地故障修复方案，演化为面向真实复杂环境的高维足式机器人控制设计原则。
